% Problem record op_df689759a5203415. This TeX file is derived from
% database/problems_json/op_df689759a5203415.json by scripts/sync-tex.mjs; edit the JSON record
% and rerun the script rather than editing this file.
\section{Sample complexity of low-rank tomography with Pauli measurements}
\label{sec:op_df689759a5203415}

\paragraph{Problem.}
How many copies of an unknown low-rank $N$-qubit state does tomography with Pauli measurements need?

Let $d=2^N$ and let $\rho$ be a density operator on $(\mathbb{C}^2)^{\otimes N}$ with $\operatorname{rank}(\rho)\leq r$, where $1\leq r\leq d$ is known to the algorithm. A Pauli measurement is labeled by a string $s\in\{X,Y,Z\}^N$. It measures every qubit $j$ in the eigenbasis of the Pauli operator $\sigma_{s_j}$ and records all $N$ outcomes $x\in\{-1,1\}^N$, so its POVM elements are
\begin{equation}
  M^{s}_{x}=\bigotimes_{j=1}^{N}\frac{I+x_j\,\sigma_{s_j}}{2},
  \qquad
  \sigma_X=X,\quad\sigma_Y=Y,\quad\sigma_Z=Z.
  \label{eq:df68-pauli-povm}
\end{equation}
An algorithm consumes one fresh copy of $\rho$ per measurement, and the string $s$ used on a copy may depend on the outcomes recorded on all earlier copies. After $n$ copies it outputs a density operator $\hat\rho$. It succeeds on $\rho$ if
\begin{equation}
  \tfrac12\lVert\hat\rho-\rho\rVert_1\leq\epsilon
  \label{eq:df68-error}
\end{equation}
holds with probability at least $2/3$. Let $n_{\mathrm{P}}(N,r,\epsilon)$ be the least $n$ for which some algorithm using only measurements of the form \eqref{eq:df68-pauli-povm} satisfies \eqref{eq:df68-error} on every state of rank at most $r$. The success probability is fixed at $2/3$. Comparisons below with bounds stated at other constant success probabilities allow a constant rescaling of $\epsilon$.

Determine how $n_{\mathrm{P}}(N,r,\epsilon)$ grows with $N$ and $r$. In exponent form, fix a sufficiently small constant $\epsilon>0$ and for $0\leq\alpha\leq1$ set
\begin{equation}
  c(\alpha):=\lim_{N\to\infty}\frac{1}{N}\log_2 n_{\mathrm{P}}\bigl(N,\lceil2^{\alpha N}\rceil,\epsilon\bigr),
  \label{eq:df68-exponent}
\end{equation}
with the limit inferior and the limit superior taken separately if the limit does not exist. The endpoints are known, $c(0)=1$ and $c(1)=\log_2 10$. For $0<\alpha<1$ the known bounds are
\begin{equation}
  \max\{1+\alpha,\ \alpha\log_2 10\}\leq c(\alpha)\leq\min\{\log_2 3+2\alpha,\ \log_2 10\},
  \label{eq:df68-known-bounds}
\end{equation}
and the lower bound improves to $\max\{1+2\alpha,\alpha\log_2 10\}$ for nonadaptive algorithms, those whose strings $s$ are all fixed before the first copy is measured. Determine $c(\alpha)$ for every $\alpha\in(0,1)$, both for the adaptive algorithms defined above and for nonadaptive algorithms. Upper and lower bounds on $n_{\mathrm{P}}(N,r,\epsilon)$ that agree up to a factor polynomial in $N$ for every $r$ would settle the question.

\subsection*{Status}

Unsolved

\subsection*{Source}

Acharya, Dharmavarapu, Liu, and Yu, in Section~1.2 of the arXiv version of their STOC~2025 paper, name the extension of Pauli-measurement tomography to low-rank states as a direction for future work and collect the bounds available for intermediate rank \sourcecite{ref:df68-acharya-stoc}{ADLY25}. The exponent \eqref{eq:df68-exponent} and the bounds \eqref{eq:df68-known-bounds} are formulated here from those bounds together with the pure-state result of Grewal, Gupta, He, Sen, and Singhal \sourcecite{ref:df68-grewal}{GGH+26} and the full-rank lower bound of Acharya, Dharmavarapu, Liu, and Yu \sourcecite{ref:df68-acharya-single-qubit}{ADLY25b}. No source states the question in this form.

\subsection*{Progress}

\begin{itemize}
  \item Full rank. With every string $s$ fixed in advance and each of the $3^N$ strings applied to an equal share of the copies, $n=O(10^N\log(1/\delta)/\epsilon^2)$ copies suffice to reach \eqref{eq:df68-error} with probability $1-\delta$ for every state \sourcecite{ref:df68-yu}{Yu20}, \sourcecite{ref:df68-acharya-stoc}{ADLY25}. Every adaptive scheme built from arbitrary single-qubit measurements, of which \eqref{eq:df68-pauli-povm} is a special case, needs $n=\Omega(10^N/(\sqrt N\,\epsilon^2))$ copies \sourcecite{ref:df68-acharya-single-qubit}{ADLY25b}. Hence $c(1)=\log_2 10$, and the full-rank gap is the factor $\sqrt N$. The lower bound appears in a preprint that has not been peer reviewed.

  \item Pure states. Given $\tilde O(2^N\log(1/\delta)/\epsilon)$ copies of a pure state, Pauli measurements with strings drawn at random in advance return a pure estimate with fidelity at least $1-\epsilon$ with probability $1-\delta$ (Theorem~1 of \sourcecite{ref:df68-grewal}{GGH+26}). For pure $\rho$ and $\hat\rho$ the trace distance equals $\sqrt{1-F}$, so $\tilde O(2^N/\epsilon^2)$ copies achieve \eqref{eq:df68-error}. Any scheme, even one measuring all copies jointly, needs $\tilde\Omega(dr/\epsilon^2)$ copies for states of rank $r$, where the tilde hides a logarithmic factor \sourcecite{ref:df68-haah}{HHJ+17}. Hence $c(0)=1$.

  \item Intermediate rank, upper bounds. Projected least squares on the outcomes of all $3^N$ Pauli measurements, each applied to an equal share of the copies, reaches \eqref{eq:df68-error} with probability $1-\delta$ once $n\geq 43\,g(d)\,r^2\epsilon^{-2}\log(d/\delta)$, with $g(d)\simeq d^{1.6}$ for Pauli basis measurements (Theorem~1 of \sourcecite{ref:df68-guta}{GKKT20}). Section~1.2 of \sourcecite{ref:df68-acharya-stoc}{ADLY25} restates this as $n=O(N\cdot 3^N r^2/\epsilon^2)$. Together with the rank-independent bound $O(10^N/\epsilon^2)$ this gives the upper bound in \eqref{eq:df68-known-bounds}. The projected least squares bound is not tight at either endpoint. It gives $N\cdot 3^N$ against $2^N$ at $r=1$ and $N\cdot 12^N$ against $10^N$ at $r=d$.

  \item Intermediate rank, lower bounds. Any nonadaptive single-copy scheme needs $\Omega(r^2 d/\epsilon^2)$ copies for states of rank $r$ (Theorem~4.11 of \sourcecite{ref:df68-lowe-nayak}{LN25}), which gives the term $1+2\alpha$ for nonadaptive algorithms. For adaptive single-copy schemes the rank-$r$ rate is unknown, since the adaptive bound $\Omega(d^3/\epsilon^2)$ of \sourcecite{ref:df68-chen}{CHL+23} covers full rank only, a gap recorded in Section~1.2 of \sourcecite{ref:df68-acharya-stoc}{ADLY25}. The term $1+\alpha$ in \eqref{eq:df68-known-bounds} therefore rests on the bound $\tilde\Omega(dr/\epsilon^2)$ for measurements entangled across copies \sourcecite{ref:df68-haah}{HHJ+17}. The term $\alpha\log_2 10$ follows from the full-rank bound of \sourcecite{ref:df68-acharya-single-qubit}{ADLY25b} by padding. It is a corollary drawn here and is not stated in that paper. For $r\geq2$ let $m=\lfloor\log_2 r\rfloor$ and embed an arbitrary $m$-qubit state $\tau$ as $\tau\otimes|0\rangle\langle0|^{\otimes(N-m)}$, a state of rank at most $r$. One copy of $\tau$ simulates one measurement \eqref{eq:df68-pauli-povm} of the embedded state, since each padded qubit returns $+1$ when measured in the $Z$ basis and an independent fair coin in the $X$ or $Y$ basis, so the transcript and every adaptive choice made from it have the correct distribution. Tracing the padded qubits out of the estimate does not increase the trace distance. Hence $n_{\mathrm{P}}(N,r,\epsilon)\geq n_{\mathrm{P}}(m,2^m,\epsilon)$. Theorem~1.1 of \sourcecite{ref:df68-acharya-single-qubit}{ADLY25b} is stated for the trace norm without the factor $\tfrac12$ and for success probability $0.9$. An algorithm for \eqref{eq:df68-error} is therefore run a constant number of times and an estimate within $2\epsilon$ of a majority of the others is selected, which raises the success probability to $0.9$ at the price of error $3\epsilon$, or $6\epsilon$ in the convention of that theorem. The hard instances of the theorem are perturbations of the maximally mixed state that must remain states, so the corollary is claimed only for $\epsilon$ below a universal constant. For such $\epsilon$ it gives $n_{\mathrm{P}}(N,r,\epsilon)=\Omega(10^m/(\sqrt m\,\epsilon^2))$, and $10^m\geq r^{\log_2 10}/10$ yields $c(\alpha)\geq\alpha\log_2 10$. This exceeds $1+\alpha$ for $\alpha>1/(\log_2 10-1)\approx0.431$ and $1+2\alpha$ for $\alpha>1/(\log_2 10-2)\approx0.756$. No other lower bound specific to Pauli measurements depends on $r$.

  \item Two-outcome Pauli measurements. When each copy is used to measure one Pauli observable, so that only the product of the outcomes in \eqref{eq:df68-pauli-povm} is kept, $O(r^2d^2\log d/\epsilon^2)$ copies suffice (Theorem~2 of \sourcecite{ref:df68-flammia}{FGLE12}), $\Omega(r^2d^2/\log d)$ copies are necessary at fixed $\epsilon$, with a constant that depends on $\epsilon$, even for adaptive sequences of observables (Theorem~6 of \sourcecite{ref:df68-flammia}{FGLE12}), and $\Omega(r^2d^2/\epsilon^2)$ copies are necessary for $\epsilon<1/8$ and nonadaptive measurements with a constant number of outcomes (Theorem~4.12 of \sourcecite{ref:df68-lowe-nayak}{LN25}). Thus, for sufficiently small fixed $\epsilon$, both adaptive and nonadaptive two-outcome Pauli tomography need $\tilde\Theta(r^2 4^N)$ copies, and for nonadaptive schemes the $\epsilon^{-2}$ dependence is also tight up to logarithmic factors. At full rank this is $16^N$ against $10^N$ for \eqref{eq:df68-pauli-povm}. The upper bound transfers to \eqref{eq:df68-pauli-povm}. The lower bounds do not.
\end{itemize}

\subsection*{References}

\begin{enumerate}[label={},leftmargin=4.8em,labelsep=0.5em]
  \footnotesize
  \item[\textup{[ADLY25]}]\label{ref:df68-acharya-stoc}
  J. Acharya, A. Dharmavarapu, Y. Liu, and N. Yu, ``Pauli measurements are not optimal for single-copy tomography,'' in \emph{Proceedings of the 57th Annual ACM Symposium on Theory of Computing (STOC 2025)}, 718--729 (2025).
\href{https://doi.org/10.1145/3717823.3718248}{doi:10.1145/3717823.3718248};
\href{https://arxiv.org/abs/2502.18170}{arXiv:2502.18170}.

  \item[\textup{[ADLY25b]}]\label{ref:df68-acharya-single-qubit}
  J. Acharya, A. Dharmavarapu, Y. Liu, and N. Yu, ``Pauli measurements are near-optimal for single-qubit tomography,'' arXiv preprint (2025).
\href{https://doi.org/10.48550/arXiv.2507.22001}{doi:10.48550/arXiv.2507.22001};
\href{https://arxiv.org/abs/2507.22001}{arXiv:2507.22001}.

  \item[\textup{[Yu20]}]\label{ref:df68-yu}
  N. Yu, ``Sample efficient tomography via Pauli measurements,'' arXiv preprint (2020).
\href{https://doi.org/10.48550/arXiv.2009.04610}{doi:10.48550/arXiv.2009.04610};
\href{https://arxiv.org/abs/2009.04610}{arXiv:2009.04610}.

  \item[\textup{[GKKT20]}]\label{ref:df68-guta}
  M. Gu\c{t}\u{a}, J. Kahn, R. Kueng, and J. A. Tropp, ``Fast state tomography with optimal error bounds,'' \emph{Journal of Physics A: Mathematical and Theoretical} \textbf{53}, 204001 (2020).
\href{https://doi.org/10.1088/1751-8121/ab8111}{doi:10.1088/1751-8121/ab8111};
\href{https://arxiv.org/abs/1809.11162}{arXiv:1809.11162}.

  \item[\textup{[GGH+26]}]\label{ref:df68-grewal}
  S. Grewal, M. Gupta, W. He, A. Sen, and M. Singhal, ``Nearly time-optimal pure state tomography with Pauli measurements,'' arXiv preprint (2026).
\href{https://doi.org/10.48550/arXiv.2601.04444}{doi:10.48550/arXiv.2601.04444};
\href{https://arxiv.org/abs/2601.04444}{arXiv:2601.04444}.

  \item[\textup{[LN25]}]\label{ref:df68-lowe-nayak}
  A. Lowe and A. Nayak, ``Lower bounds for learning quantum states with single-copy measurements,'' \emph{ACM Transactions on Computation Theory} \textbf{17}(1), Article 7 (2025).
\href{https://doi.org/10.1145/3717450}{doi:10.1145/3717450};
\href{https://arxiv.org/abs/2207.14438}{arXiv:2207.14438}.

  \item[\textup{[HHJ+17]}]\label{ref:df68-haah}
  J. Haah, A. W. Harrow, Z. Ji, X. Wu, and N. Yu, ``Sample-optimal tomography of quantum states,'' \emph{IEEE Transactions on Information Theory} \textbf{63}(9), 5628--5641 (2017).
\href{https://doi.org/10.1109/TIT.2017.2719044}{doi:10.1109/TIT.2017.2719044};
\href{https://arxiv.org/abs/1508.01797}{arXiv:1508.01797}.

  \item[\textup{[CHL+23]}]\label{ref:df68-chen}
  S. Chen, B. Huang, J. Li, A. Liu, and M. Sellke, ``When does adaptivity help for quantum state learning?,'' in \emph{2023 IEEE 64th Annual Symposium on Foundations of Computer Science (FOCS)}, 391--404 (2023).
\href{https://doi.org/10.1109/FOCS57990.2023.00029}{doi:10.1109/FOCS57990.2023.00029};
\href{https://arxiv.org/abs/2206.05265}{arXiv:2206.05265}.

  \item[\textup{[FGLE12]}]\label{ref:df68-flammia}
  S. T. Flammia, D. Gross, Y.-K. Liu, and J. Eisert, ``Quantum tomography via compressed sensing: error bounds, sample complexity and efficient estimators,'' \emph{New Journal of Physics} \textbf{14}, 095022 (2012).
\href{https://doi.org/10.1088/1367-2630/14/9/095022}{doi:10.1088/1367-2630/14/9/095022};
\href{https://arxiv.org/abs/1205.2300}{arXiv:1205.2300}.
\end{enumerate}

\subsection*{Comment}

The open range is $0<\alpha<1$. There is evidence that neither side of \eqref{eq:df68-known-bounds} is tight. The projected least squares bound misses the truth by an exponential factor at both endpoints. The lower bound $1+2\alpha$ holds for every nonadaptive single-copy scheme, whereas at $\alpha=1$ Pauli measurements already lose the factor $(10/8)^N$ against the best single-copy scheme, and the padding bound $\alpha\log_2 10$ treats the $N-m$ padded qubits as free. Whether $c$ is convex or piecewise linear is not known. To give concrete numbers, at $r=\sqrt d$ the nonadaptive lower bound is $4^N/\epsilon^2$ copies, the upper bound is $N\cdot 6^N/\epsilon^2$ copies, and the adaptive lower bound is $10^{N/2}/(\sqrt N\,\epsilon^2)$ copies.

Nearby questions remain open as well. In the adaptive model the lower bound for $\alpha<0.431$ is only $1+\alpha$, because no adaptive single-copy lower bound is known beyond full rank. At full rank the factor $\sqrt N$ between $10^N/\epsilon^2$ and $10^N/(\sqrt N\,\epsilon^2)$ remains, and \sourcecite{ref:df68-acharya-single-qubit}{ADLY25b} asks whether $10^N$ is the sample complexity of Pauli measurements. The same question with infidelity in place of \eqref{eq:df68-error} is listed with the low-rank extension in \sourcecite{ref:df68-acharya-stoc}{ADLY25}. The two-outcome model of the last Progress item is settled up to logarithmic factors at fixed $\epsilon$ and is a different problem from \eqref{eq:df68-pauli-povm}.

\subsection*{Field}

Quantum metrology

\subsection*{Topic}

Quantum estimation

\subsection*{ID}

\texttt{op\_df689759a5203415}
